% NOTE: single paragraph, plain language, no citations. Target 200--250 words.
% The two \pending sentences must be rewritten with observed effect sizes after
% the authenticated four-block campaign; do not soften the pending markers.
Scour --- the erosion of soil around bridge foundations during floods --- is a
leading cause of bridge failure, yet it is difficult to inspect exactly when it
matters most. Drive-by monitoring offers an attractive alternative: an
in-service train records its own vibration while crossing the bridge, and that
response carries information about support condition. Previous studies
established feasibility but leave pre-deployment choices unresolved:
whether to process raw or compressed signals, which neural components help,
and which response channels are informative. This study addresses
those questions with a prospectively specified, source-locked experiment based
on a two-dimensional train--track--bridge simulation of a two-span benchmark
and a four-span geometry. Four blocks separate scour
from a multi-damage task that activates nominal bearing fixity and a local
flexural-rigidity-loss surrogate. One fixed rail-profile realization is shared
throughout; speed, temperature, and vehicle properties vary between passages,
while track damage, wheel polygonization, and suspension damage are deferred.
A 16-cell factorial compares raw and Piecewise Aggregate Approximation inputs,
positional encoding, recurrence, and fixed-width multi-rate pooling. A sealed,
state-grouped protocol then screens all eight single channels descriptively and
all 28 two-channel pairs for four retained pipelines. The eight modeled inputs
include two constrained-wheelset acceleration proxies and therefore represent
responses, not a validated eight-sensor package.
\pending{State the authenticated RAW/PAA factorial and selected-pair outcomes,
with effect sizes and explicitly descriptive finite-design uncertainty.}
\pending{State the authenticated block-local scour, multi-damage, and
multi-pier results without importing pilot or pre-freeze values.} The resulting
workflow provides an auditable template for simulation-based design decisions
before field validation.
